\documentclass[11pt]{article}
\usepackage[T1]{fontenc}
\usepackage{lmodern,microtype,amsmath,amssymb,amsthm,needspace}
\usepackage[margin=1in]{geometry}
\usepackage[colorlinks=true,linkcolor=blue,urlcolor=blue,citecolor=blue]{hyperref}
\emergencystretch=2em
\hypersetup{pdftitle={A smooth Reshetnyak-harmonic sphere with unbounded distortion},pdfauthor={Alec Kriebel}}
\newtheorem{theorem}{Theorem}
\newtheorem{lemma}[theorem]{Lemma}
\newcommand{\Ep}{E_+^2}
\newcommand{\EK}{E_{\rm KS}}
\newcommand{\R}{\mathbb R}
\newcommand{\tr}{\operatorname{tr}}
\newcommand{\op}{\operatorname{op}}
\title{A smooth Reshetnyak-harmonic sphere\\with unbounded distortion}
\author{Alec Kriebel\\\small Independent researcher\\\small\href{https://orcid.org/0009-0001-9320-500X}{ORCID: 0009-0001-9320-500X}}
\date{13 October 2026 \quad Version 1.0\\\small Unrefereed research note}
\begin{document}
\maketitle
\begin{abstract}
We construct a compact smooth Riemannian sphere and a smooth homeomorphism
from the round sphere which locally minimizes Reshetnyak energy against every
same-trace Sobolev competitor, but has essentially unbounded infinitesimal
distortion. This gives a negative answer for the Reshetnyak member of the
general-energy harmonic-sphere question posed in the inspected 2025 texts of
Meier--Vikman--Wenger. Scalar calibration on a flat belt and an exact two-form
calibration on conformal caps prove the full local comparison. The map does
not minimize energy in its homotopy class and is not Korevaar--Schoen harmonic.
Separately, an angular-stress argument gives the nonsharp bound
$4+\sqrt{17}$ for locally Korevaar--Schoen-minimizing spheres into complete
metric spaces; the sharper classical SL(2)-minimal bound also applies.
Classical ingredients are credited; first priority is not claimed.
\end{abstract}

\section{Question, energies, and exact conclusions}
Question~1 of Meier's Oberwolfach contribution \cite[p.~1625]{owr}, also
Question~1.9 of Meier--Vikman--Wenger \cite{mvw}, asks whether harmonic spheres
are infinitesimally quasiconformal in its stated metric setting. The energy is
fixed from a class including Reshetnyak and Korevaar--Schoen energies
\cite[Definition~3.1]{mvw}. We use the precise neighborhoodwise definition:
each source point has an open neighborhood $U$ such that, for every Lipschitz
domain $\Omega\Subset U$ and every $v\in W^{1,2}(\Omega,X)$ with the trace of
$u|_\Omega$, the chosen energy of $u|_\Omega$ is at most that of $v$
\cite[Definition~6.1]{mvw}. Competitors may leave the reference image and need
not be continuous or homotopic to the reference map.

For a smooth Riemannian target, the energies in conformal source coordinates are
\[
 \Ep(v;\Omega)=\int_\Omega\|dv\|_{\op}^2,
 \qquad \EK(v;\Omega)=\int_\Omega(\sigma_1(dv)^2+\sigma_2(dv)^2).
\]
The conformal source factor cancels. For a metric differential seminorm $s$,
let $M(s)$ and $m(s)$ be its maximum and minimum on the Euclidean unit circle.
Infinitesimal quasiconformality with constant $Q$ means
$M(s)\le Qm(s)$ almost everywhere; the zero seminorm satisfies this inequality.

\begin{theorem}\label{thm:negative}
There is a compact smooth Riemannian surface $X$, diffeomorphic to $S^2$,
and a smooth homeomorphism $u:S^2\to X$ such that $u$ is locally minimizing
for $\Ep$ in the preceding full Sobolev sense, while no finite $Q$ bounds its
infinitesimal distortion. The target is geodesic, admits a local quadratic
isoperimetric inequality, and every sufficiently small-diameter continuous
sphere in $X$ is null-homotopic.
\end{theorem}

\begin{theorem}\label{thm:positive}
Let $X$ be complete and $u\in W^{1,2}(S^2,X)$ be locally minimizing for
Korevaar--Schoen energy in the standard metric-Sobolev framework of
\cite{mvw}. Then $u$ is infinitesimally $(4+\sqrt{17})$-quasiconformal.
No target curvature, topology, or isoperimetric hypothesis is needed.
\end{theorem}

Theorem~\ref{thm:negative} answers the inspected general-energy question
negatively. Theorem~\ref{thm:positive} addresses a separate fixed-energy
restriction. Energy comparability does not identify their exact minimizers.
The source binding is to the byte-checked arXiv v1 and the Oberwolfach report.
The revised 2026 journal article exists, but its mathematical text was not
accessible to this review; unchanged final wording or openness through 2026
is not asserted.

\section{Two local comparisons with unrestricted competitors}
\begin{lemma}\label{lem:flat}
Suppose a target patch is Euclidean with coordinates $(s,y)$, and a smaller
set $V$ has $y$ 1-Lipschitz for the whole target distance. In source coordinates
$(t,x)$, a map $u(t,x)=(h(t),x)$ into $V$ with $|h'|\le1$ minimizes $\Ep$
on every relatively compact Lipschitz subdomain against all same-trace Sobolev
competitors in the whole target.
\end{lemma}
\begin{proof}
The McShane extension \cite{mcshane}
\[
 F(q)=\inf_{p\in V}\{y(p)+d_X(p,q)\}
\]
is globally 1-Lipschitz and agrees with $y$ on $V$; it is finite in the compact
target used below. For any competitor $v$, scalar Sobolev composition gives
$f=F\circ v\in W^{1,2}(\Omega)$, with trace $x$ and
$\int_\Omega|\nabla f|^2\le\Ep(v;\Omega)$. Put
$w=f-x\in W^{1,2}_0(\Omega)$. Integration of its weak $x$-derivative is zero,
by compactly supported approximation. Therefore
\[
 \Ep(v;\Omega)\ge\int_\Omega|\nabla f|^2
 =|\Omega|+\int_\Omega|\nabla w|^2\ge|\Omega|
 =\Ep(u;\Omega).
\]
There is no restriction on the competitor's image. In a smooth target, a
sufficiently small strongly geodesically convex ball inside a flat coordinate
patch supplies the ambient-distance premise: its minimizing geodesics remain
in the patch, so coordinate differences are bounded by their lengths.
\end{proof}

\begin{lemma}\label{lem:cap}
A smooth orientation-preserving conformal map into an oriented compact smooth
Riemannian sphere locally minimizes $\Ep$ against every same-trace
$W^{1,2}$ competitor.
\end{lemma}
\begin{proof}
Write $A=\operatorname{Area}(X)$. Shrink a source neighborhood $U$ until its
compact image supports a smooth cutoff $0\le\eta\le1$ which equals one near
that image and has $0<a:=\int_X\eta\,dA<A/2$. Set
\[
 b=\frac a{A-a}<1,\qquad \omega=((1+b)\eta-b)\,dA.
\]
This two-form has comass at most one, agrees with the area form near $u(U)$,
and integrates to zero on $X$. Top-degree de Rham cohomology gives a global
smooth one-form $\alpha$ with $d\alpha=\omega$.
The standard calibration comparison \cite{hl} gives, pointwise for every
competitor,
\[
 v^*\omega\le |\operatorname{Jac}v|\,dA_{\rm source}
 \le\|dv\|_{\op}^2\,dA_{\rm source};
\]
both inequalities are equalities for the reference conformal map. We justify
the needed identity $\int_\Omega v^*\omega=\int_\Omega u^*\omega$ at the
critical Sobolev regularity, rather than assuming smooth target-valued density.

A finite atlas and partition of unity write $\alpha=\sum_j f_j\,dg_j$ with
globally smooth real functions on $X$. Coordinate functions are extended with
cutoffs equal to one near the corresponding coefficient supports. For bounded
scalar $F,G\in W^{1,2}_{\rm loc}$, mollification proves
$d(F\,dG)=dF\wedge dG$ in distributions: uniformly bounded approximants for
$F$ give $L^2$ convergence of the one-forms by dominated convergence and
strong $L^2$ gradient convergence; products of the gradients converge in $L^1$.
Applying scalar composition and the chain rule thus gives
$d(v^*\alpha)=v^*\omega$.

Matching traces glue $v$ on $\Omega$ to $u$ on $U\setminus\Omega$ as a
Sobolev map $\widetilde v$. This can also be checked in a smooth embedding of
the compact target. The $L^2$ one-form
$\beta=\widetilde v^*\alpha-u^*\alpha$ is supported in
$\overline\Omega\Subset U$. Pair its distributional derivative with a smooth
cutoff equal to one there. The cutoff derivative misses its support, so
\[
 0=\int_U d\beta=\int_\Omega(v^*\omega-u^*\omega),
\]
where the middle integral denotes that cutoff pairing. Integrating the
calibration inequalities proves the assertion. Extra bubbles, orientation
reversal, and target excursions do not alter this exact-form identity.
\end{proof}

\section{The smooth compact counterexample}
Use the cylinder $\R\times(\R/2\pi\mathbb Z)$ with coordinates $(s,\theta)$
and add its two ends as sphere poles. Define
\[
 B(r)=\begin{cases}e^{-1/r}&r>0,\\0&r\le0,\end{cases}
 \quad H(r)=\frac{B(r)}{B(r)+B(1-r)},\quad
 \rho(s)=e^{-|s|H(|s|-2)},\quad
 g_X=\rho(s)^2(ds^2+d\theta^2).
\]
The denominator never vanishes. The metric is flat for $|s|\le2$. At the
negative end $z=e^{s+i\theta}$ makes it exactly $|dz|^2$ for $s\le-3$;
at the positive end $\zeta=e^{-s-i\theta}$ does the same for $s\ge3$.
It therefore extends smoothly and positively over both poles. The apparent
absolute value is harmless since its coefficient vanishes near zero.

Compact smooth Riemannian geometry gives completeness and geodesicity, and a
uniform contractible normal-ball radius. Every sufficiently small-diameter
sphere image is contained in such a ball. To verify local quadratic isoperimetry,
choose uniformly controlled nested charts with convex Euclidean images.
A sufficiently short Lipschitz loop lies in a smaller chart. Its Euclidean
image of length at most $L\ell$ can be coned to one of its points within
distance $L\ell/2$, giving cone area at most $L^2\ell^2/4$. Returning through
the inverse chart gives area at most $L^4\ell^2/4$. The radial cone is Sobolev
and retains the prescribed Lipschitz boundary parametrization. Compactness
makes the scale and constants uniform.

On the round source sphere use cylinder coordinates $(t,\theta)$, with metric
$g_0=\operatorname{sech}^2t\,(dt^2+d\theta^2)$. Put
\[
 \eta_0(t)=1-H((16t^2-1)/3),\quad
 a(t)=\eta_0(t)t^2+1-\eta_0(t),\quad h(t)=\int_0^t a(r)\,dr.
\]
Here $a$ is smooth and even, $0\le a\le1$, vanishes only at zero, equals $t^2$
for $|t|\le1/4$, and equals one for $|t|\ge1/2$. Thus $h$ is a smooth,
strictly increasing odd homeomorphism of $\R$, with
\begin{align*}
 h(t)&=t^3/3\quad(|t|\le1/4),\\
 h(t)&=t-c\quad(t\ge1/2),\qquad h(t)=t+c\quad(t\le-1/2),\\
 c&=\int_0^{1/2}(1-a(r))\,dr\in(0,1/2).
\end{align*}
Define $u(t,\theta)=(h(t),\theta)$ and send the poles to the corresponding
target poles. Its pole expressions are multiplication by $e^c$ in the
respective complex coordinates, so it is a globally smooth homeomorphism.
The inverse need not be smooth at the equator; that is not a hypothesis.
Smoothness on the compact source implies $u\in W^{1,2}(S^2,X)$.

For $|t|<1$, the image is in the flat target belt since $|h(t)|\le|t|$.
Around each such point choose a smaller source neighborhood mapping into a
strongly convex target ball, and lift the angular coordinate to a real $x$.
Lemma~\ref{lem:flat} proves every required fixed-trace comparison there.
For $|t|>1/2$ and at the poles, the map is orientation-preserving conformal;
Lemma~\ref{lem:cap} proves every required comparison. These open regimes
overlap and cover the sphere, including both seam circles.

For $0<|t|<1/4$, the singular values before the common source conformal
rescaling are $1$ and $t^2$. Their ratio is $t^{-2}$. Given finite $Q\ge1$,
choose $0<\delta<\min\{1/4,1/\sqrt Q\}$. The band $0<|t|<\delta$ has
positive round area $4\pi\tanh\delta$ and ratio greater than $Q$.
This proves essential unboundedness and Theorem~\ref{thm:negative}.

For completeness, the maximum-square seminorm energy is admissible: continuity,
monotonicity, degree-two homogeneity, and rotation invariance are immediate.
Its unit sublevel consists of uniformly bounded 1-Lipschitz seminorms on the
unit circle, compact by Arzel\`a--Ascoli. Quasiconvexity follows from convexity
of the squared operator norm: a smooth fixed-affine-boundary map into a normed
space has average derivative equal to that affine derivative, and Jensen's
inequality gives the energy comparison. No immersion restriction is needed.

\subsection*{Two distinctions that prevent false contradictions}
The conformal coordinate identity $f(t,\theta)=(t,\theta)$ is homotopic to $u$
and has $\Ep(f)=\operatorname{Area}(X)$. Changing variables $s=h(t)$ gives
\[
 \Ep(u)-\operatorname{Area}(X)
 =2\pi\int_{-1/2}^{1/2}(1-a(t))\,dt=4\pi c>0.
\]
Thus $u$ is not a global homotopy minimizer. In a flat rectangle with
$0<t<1/4$, perturb its first coordinate by $\varepsilon\psi$, where
$\psi\ge0$ is smooth, nonzero, and compactly supported. The KS first variation is
\[
 2\int h'\partial_t\psi=-2\int h''\psi=-4\int t\psi<0.
\]
Small perturbations keep the image flat and retain the trace, so $u$ is not
KS-harmonic. Nonsmooth operator-norm energy does not supply the unique linear
stress used next; local supporting stresses cannot simply be patched globally.

\section{The separate Korevaar--Schoen theorem}
The metric-Sobolev approximate differential and domain-diffeomorphism chain
rule are standard inputs \cite{ks,mvw,lw}. In a conformal coordinate chart,
$s_z$ denotes the approximate metric differential of $u$ at $z$ on the
coordinate tangent plane $\R^2$. Normalize KS energy by
$I(s)=\pi^{-1}\int_{S^1}s(v)^2\,dv$. For $A\in GL^+(2)$, set
$F_s(A)=I(s\circ A)/\det A$. The angular substitution $w=Av/|Av|$, whose
Jacobian is $\det A/|Av|^2$, yields
\begin{equation}\label{eq:angular}
 F_s(A)=\frac1{\pi(\det A)^2}\int_{S^1}
             \frac{s(w)^2}{|A^{-1}w|^4}\,dw.
\end{equation}
This is smooth in $A$ near the identity even for a nonsmooth or degenerate
seminorm. Each fixed-order matrix derivative is bounded by $C M(s)^2$.
Its first derivative is $DF_s(I)[D]=T(s):D$, where
\[
 T(s)=\frac1\pi\int_{S^1}s(w)^2(4w\otimes w-2I)\,dw.
\]
The stress is symmetric and tracefree. For a smooth vector field $V$ supported
inside a minimizing source chart, use $u\circ\phi_\varepsilon$ with
$\phi_\varepsilon=\mathrm{id}+\varepsilon V$. On a smaller domain enclosing
its support this is an admissible same-trace competitor for either sign of
$\varepsilon$. After the change of variables, the differential seminorm
is held at the fixed new point and the matrix is
$D\phi_\varepsilon\circ\phi_\varepsilon^{-1}$. The $CM(s)^2$ bound is
integrable. Dominated differentiation gives
\[
 \int T(s_z):DV(z)\,dz=0.
\]
Write $T=\left(\begin{smallmatrix}a&b\\b&-a\end{smallmatrix}\right)$.
Its distributional divergence equations are
$a_x+b_y=0$, $b_x-a_y=0$, so $q=a-ib$ is holomorphic by the distributional
Weyl lemma. Scaling and rotation give
$T(\lambda s)=\lambda^2T(s)$ and $T(s\circ R)=R^TT(s)R$.
Under a holomorphic coordinate change $z=z(w)$,
$q_w=(dz/dw)^2q_z$. Thus $q_z\,dz^2$ is a global holomorphic quadratic
differential. Each pole has its own minimizing chart and weak equation;
no unproved puncture-removal claim is used. On the sphere this differential
vanishes: its coefficient is entire in a plane chart and is $O(|z|^{-4})$
at infinity. Hence $T(s_z)=0$ almost everywhere.

For a nonzero balanced seminorm normalize $M(s)=1$. Choose $e_2$ with
$s(e_2)=m$, and put $L=s(e_1)$ for its perpendicular unit vector.
The triangle inequality gives $L^2\ge1-m^2$. For $r(v)=L|v_1|$,
the reverse triangle inequality gives $|s(v)-r(v)|\le m$ on the unit circle.
Both are at most one, so $|s(v)^2-r(v)^2|\le2m$ and
\[
 \|T(s)-T(r)\|_{\op}\le8m,\qquad T(r)=L^2\operatorname{diag}(1,-1).
\]
Consequently $1-m^2\le L^2\le8m$, and $m\ge\sqrt{17}-4$.
Undoing normalization proves $M/m\le4+\sqrt{17}$; zero seminorms satisfy
the inequality trivially. This proves Theorem~\ref{thm:positive}.

For inner-product seminorms, $T(s)=2G-\tr(G)I$ and this recovers classical
weak conformality \cite[Lemma~1.5]{su}. In general, fourfold-symmetric
non-Euclidean norms can be balanced. Bounded distortion, not roundness, is the
conclusion. The older, sharper KS bound also applies here. Indeed,
$J(A)=I(s\circ A)$ is convex in the matrix $A$, and the angular identity gives
$DJ(I)[D]=T(s):D+I(s)\tr D$. Once $T(s)=0$, for any positive symmetric
$S$ with $\det S=1$ convexity and $\tr S\ge2$ imply
\[
 J(S)\ge J(I)+I(s)(\tr S-2)\ge J(I).
\]
Every $A\in SL(2)$ has polar decomposition $A=SR$, with $R$ a rotation;
circular invariance gives $J(A)=J(S)$. Thus the balanced seminorm is
SL(2)-minimal. The pointwise Lytchak--Wenger bound
$2\sqrt2+\sqrt6$ applies \cite[Lemma~3.1 and \S3.3]{lw}.
We retain the elementary $4+\sqrt{17}$ proof for a self-contained bound;
no new, improved, or optimal distortion constant is claimed.

\section*{Attribution, verification, and review status}
The question and homotopy-minimizer theory belong to Meier--Vikman--Wenger.
We credit Lytchak--Wenger's general energy/area framework, Korevaar--Schoen's
Sobolev and NPC tensor theory, the classical Hopf/Sacks--Uhlenbeck method,
McShane extension, and standard form/comass calibration. The compact belt/cap
construction and the complete arguments above are submitted for review.
The accompanying bounded priority audit did not authenticate an earlier exact
match, but first bibliographic priority is unresolved. Final publisher wording,
inaccessible older texts, and unread or unindexed later literature remain
historical coverage gaps. No first-priority or current-openness claim is made.

The source and verification archive preserves the original candidate and
independent analytic, computational, source/model, and priority reviews.
The historical author has 257 exact algebra controls. The source-free replay
executes 673 of a historical 675-check mixed audit, omitting its two PDF-input
identity checks; another mixed audit is reproduced within its stated numerical
tolerances. Fresh independent controls and negative tests are included.
Numerical quadrature is diagnostic, not a rigorous error enclosure. These
finite checks supplement the universal written proofs and do not formally
certify Sobolev analysis or historical priority.

AI tools were used extensively in development, computation, verification,
adversarial review, and writing. This note is unrefereed and has not received
independent human or journal peer review. AI-assisted independent review is
not human peer review or proof-assistant formalization. Authored materials are
released under CC BY 4.0; cited third-party documents are not included.

\begin{thebibliography}{9}
\bibitem{mvw} D.~Meier, N.~Vikman, S.~Wenger,
\emph{Energy minimizing harmonic $2$-spheres in metric spaces},
arXiv:2503.08553v1, 11 March 2025.
\url{https://arxiv.org/pdf/2503.08553v1}.
\bibitem{owr} D.~Meier (joint work with N.~Vikman and S.~Wenger),
\emph{Energy minimizing harmonic $2$-spheres in metric spaces},
Oberwolfach Report 30/2025, pp.~1624--1626, Question 1 on p.~1625.
\url{https://doi.org/10.4171/owr/2025/30}.
\bibitem{lw} A.~Lytchak, S.~Wenger,
\emph{Energy and area minimizers in metric spaces},
Adv. Calc. Var.~10 (2017), 407--421.
\url{https://doi.org/10.1515/acv-2015-0027}.
Relevant manuscript: 11 September 2015, Lemmas 3.1/4.1 and \S3.3,
\url{https://www.mi.uni-koeln.de/~alytchak/preprints/enarea.pdf}.
\bibitem{ks} N.~J.~Korevaar, R.~M.~Schoen,
\emph{Sobolev spaces and harmonic maps for metric space targets},
Comm. Anal. Geom.~1 (1993), 561--659.
\url{https://doi.org/10.4310/cag.1993.v1.n4.a4}.
\bibitem{su} J.~Sacks, K.~Uhlenbeck,
\emph{The existence of minimal immersions of $2$-spheres},
Ann. of Math.~113 (1981), 1--24.
\url{https://doi.org/10.2307/1971131}.
\bibitem{mcshane} E.~J.~McShane, \emph{Extension of range of functions},
Bull. Amer. Math. Soc.~40 (1934), 837--842.
\url{https://doi.org/10.1090/S0002-9904-1934-05978-0}.
\bibitem{hl} R.~Harvey, H.~B.~Lawson, Jr., \emph{Calibrated geometries},
Acta Math.~148 (1982), 47--157.
\url{https://doi.org/10.1007/BF02392726}.
\end{thebibliography}
\end{document}
